\chapter{Results}

This section presents the empirical findings of the analysis.

\section{Descriptive Statistics}

Table~\ref{tab:descriptive} summarizes the distribution of malaria infection rates, ITN use, and key covariates.

\begin{table}[htbp]
	\centering
	\caption{Descriptive Statistics by Survey Year}
	\label{tab:descriptive}
	\begin{tabular}{@{}lccc@{}}
		\toprule
		\textbf{Characteristic} & \textbf{KMIS 2015} & \textbf{KMIS 2020} & \textbf{Combined} \\
		\midrule
		Number of children & 3,286 & 3,485 & 6,771 \\
		Number of clusters & 245 & 296 & 528 \\
		Malaria prevalence (\%) & 5.1 & 4.5 & 4.8 \\
		ITN use rate (\%) & 58.0 & 50.5 & 54.1 \\
		Urban population (\%) & 30.5 & 27.5 & 29.0 \\
		\bottomrule
	\end{tabular}
\end{table}

Malaria prevalence declined from 5.1\% in 2015 to 4.5\% in 2020, while ITN use also declined from 58.0\% to 50.5\%. Figure~\ref{fig:malaria_by_wealth} shows a clear gradient: the poorest children have approximately seven times higher malaria risk than the richest children.

\begin{figure}[htbp]
	\centering
	\includegraphics[width=0.8\textwidth]{Figures/malaria_by_wealth.png}
	\caption{Malaria Prevalence by Wealth Quintile (KMIS 2015 vs 2020)}
	\label{fig:malaria_by_wealth}
\end{figure}

\subsection{Positivity Check}

Following \citet{westreich2010invited}, we assess the positivity assumption using propensity score overlap plots. Figure~\ref{fig:overlap} shows substantial overlap between treated (PS range: 0.066-0.986) and control (PS range: 0.019-0.936) groups, indicating no violations.

\begin{figure}[htbp]
	\centering
	\includegraphics[width=0.8\textwidth]{Figures/ps_overlap_glmm_Combined.png}
	\caption{Propensity Score Overlap Plot (GLMM-Based, Combined Sample)}
	\label{fig:overlap}
\end{figure}

\section{Baseline Regression Results}

Table~\ref{tab:baseline} presents the baseline regression results.

\begin{table}[htbp]
	\centering
	\caption{Baseline Logistic Regression Results (Odds Ratios)}
	\label{tab:baseline}
	\begin{tabular}{@{}lccc@{}}
		\toprule
		\textbf{Model} & \textbf{KMIS 2015} & \textbf{KMIS 2020} & \textbf{Combined} \\
		\midrule
		Unadjusted & 1.105 (0.710-1.721) & 1.162 (0.760-1.778) & 1.178 (0.860-1.614) \\
		Demographic + Socioeconomic & 1.279 (0.792-2.065) & 1.551 (0.994-2.420) & 1.465 (1.051-2.042) \\
		Fully Adjusted & \textbf{0.592 (0.360-0.973)} & 0.850 (0.560-1.291) & 0.765 (0.551-1.062) \\
		\bottomrule
	\end{tabular}
\end{table}



\section{GLMM Results}

Table~\ref{tab:glmm} presents GLMM results.

\begin{table}[htbp]
	\centering
	\caption{GLMM Results for ITN Use on Malaria}
	\label{tab:glmm}
	\begin{tabular}{@{}lcccc@{}}
		\toprule
		\textbf{Survey} & \textbf{OR} & \textbf{95\% CI} & \textbf{p-value} & \textbf{ICC} \\
		\midrule
		KMIS 2015 & 0.509 & (0.322-0.805) & 0.004 & 0.429 \\
		KMIS 2020 & 0.556 & (0.375-0.824) & 0.004 & 0.414 \\
		Combined & 0.535 & (0.397-0.721) & $<$0.001 & 0.471 \\
		\bottomrule
	\end{tabular}
\end{table}

The ICC values of 41-47\% indicate substantial clustering \citep{killip2004intracluster}. The GLMM estimates show stronger protective effects than fully adjusted logistic regression, as expected for cluster-specific effects \citep{neuhaus1991comparison}.

\section{Double Machine Learning Estimates}

Table~\ref{tab:dml} presents DML estimates using Random Forest with 5-fold cross-fitting.

\begin{table}[htbp]
	\centering
	\caption{Double Machine Learning Results (Random Forest, 5 folds)}
	\label{tab:dml}
	\begin{tabular}{@{}lccc@{}}
		\toprule
		\textbf{Survey} & \textbf{OR} & \textbf{95\% CI} & \textbf{p-value} \\
		\midrule
		KMIS 2015 & 0.973 & (0.948-0.998) & 0.032 \\
		KMIS 2020 & 0.989 & (0.966-1.013) & 0.378 \\
		Combined & 0.981 & (0.963-0.999) & 0.041 \\
		\bottomrule
	\end{tabular}
\end{table}

Table~\ref{tab:dml_xgb} presents DML estimates using XGBoost.

\begin{table}[htbp]
	\centering
	\caption{Double Machine Learning Results (XGBoost, 5 folds)}
	\label{tab:dml_xgb}
	\begin{tabular}{@{}lccc@{}}
		\toprule
		\textbf{Survey} & \textbf{OR} & \textbf{95\% CI} & \textbf{p-value} \\
		\midrule
		KMIS 2015 & 0.978 & (0.961-0.996) & 0.018 \\
		KMIS 2020 & 0.976 & (0.961-0.991) & 0.002 \\
		Combined & 0.975 & (0.963-0.986) & $<$0.001 \\
		\bottomrule
	\end{tabular}
\end{table}

The DML estimates indicate a 2-4\% reduction in malaria odds at the population level. Scaled to Kenya's population of approximately 5 million children under five, this translates to approximately 125,000 malaria cases averted annually.

\subsection{Comparison with Existing Evidence}


\section{Heterogeneity Analysis}

Table~\ref{tab:heterogeneity} presents subgroup-specific treatment effects.

\begin{table}[htbp]
	\centering
	\caption{Heterogeneous Treatment Effects by Subgroup (DML)}
	\label{tab:heterogeneity}
	\begin{tabular}{@{}llccc@{}}
		\toprule
		\textbf{Subgroup} & \textbf{Category} & \textbf{N} & \textbf{OR} & \textbf{p-value} \\
		\midrule
		Wealth & Poorer & 1,448 & 0.928 & 0.002 \\
		Residence & Urban & 2,240 & 0.974 & 0.003 \\
		Age & 36-59 months & 3,075 & 0.957 & $<$0.001 \\
		Rainfall & Medium & 2,227 & 0.954 & $<$0.001 \\
		Rainfall & High & 2,156 & 0.945 & 0.001 \\
		\bottomrule
	\end{tabular}
\end{table}

The strongest protective effects are observed in the poorer wealth quintile (OR = 0.928), urban areas (OR = 0.974), older children (OR = 0.957), and high-rainfall areas (OR = 0.945).

\section{Robustness Checks}

\subsection{Alternative Machine Learning Algorithms}

All five learner configurations produced odds ratios between 0.961 and 0.982, all statistically significant, demonstrating robustness to different machine learning specifications.

\subsection{Alternative Number of Cross-Fitting Folds}

Estimates were highly stable across 5-fold and 10-fold cross-fitting, with maximum difference of 0.012.

\subsection{Alternative Covariate Specifications}

The protective effect remained consistent (OR between 0.96 and 0.98) across all covariate specifications.

Figure~\ref{fig:robustness} presents a forest plot comparing robustness across learners.

\begin{figure}[htbp]
	\centering
	\includegraphics[width=0.8\textwidth]{Figures/week6_robustness_comparison.png}
	\caption{Robustness of DML Estimates to Different Learners}
	\label{fig:robustness}
\end{figure}